\documentclass[12pt,hyperref,a4paper,UTF8,fontset=none]{ctexart}

\usepackage{zjureport}
\usepackage{enumitem}
\usepackage{listings}
\usepackage{xcolor}
\usepackage{booktabs}
\usepackage{longtable}
\usepackage{graphicx}
\usepackage{float}
\usepackage{amsmath}
\usepackage{array}
\usepackage{caption}
\usepackage{tabularx}
\usepackage{setspace}

\zjuheader{网络安全原理实践：实验二 ARP 缓存中毒攻击}
\zjucourse{网络安全原理实践}
\zjureporttitle{实验二：ARP 缓存中毒攻击}
\zjuname{周子为}
\zjuemail{zzw4257@gmail.com}
\zjucollege{竺可桢学院 / 计算机科学与技术学院}
\zjudepartment{混合班 / 信息安全}
\zjumajor{信息安全}
\zjustuid{3230106267}
\zjuinstructor{杜文亮教授}
\zjureportdate{2026年3月15日}

\setlength{\headheight}{15pt}
\setstretch{1.22}
\setlength{\parskip}{0.35em}
\hypersetup{
    colorlinks=true,
    linkcolor=black,
    filecolor=black,
    urlcolor=blue,
    citecolor=black,
}
\captionsetup{font=small, labelfont=bf}
\setlist[itemize]{leftmargin=2em, itemsep=0.2em, topsep=0.3em}
\setlist[enumerate]{leftmargin=2em, itemsep=0.2em, topsep=0.3em}
\graphicspath{{figures/}{../evidence/}{../evidence/manual/}}

\definecolor{codegray}{gray}{0.96}
\definecolor{codepurple}{rgb}{0.58,0,0.82}
\definecolor{codeblue}{rgb}{0,0,0.75}
\definecolor{codegreen}{rgb}{0,0.55,0}

\lstset{
    backgroundcolor=\color{codegray},
    commentstyle=\color{codegreen},
    keywordstyle=\color{codeblue},
    numberstyle=\tiny\color{gray},
    stringstyle=\color{codepurple},
    basicstyle=\ttfamily\small,
    breakatwhitespace=false,
    breaklines=true,
    columns=fullflexible,
    captionpos=b,
    keepspaces=true,
    numbers=left,
    numbersep=5pt,
    showspaces=false,
    showstringspaces=false,
    showtabs=false,
    frame=single,
    rulecolor=\color{black},
    tabsize=2
}

\newcommand{\manualevidencepath}[1]{../evidence/manual/#1}
\newcommand{\autoevidencepath}[1]{../evidence/#1}

\newcommand{\placeholderpanel}[2]{%
    \fbox{\begin{minipage}[c][0.18\textheight][c]{0.47\textwidth}
        \centering \textbf{截图占位}\\[0.5em]
        \small 建议文件名：\path{#1}\\
        \small #2
    \end{minipage}}%
}

\newcommand{\manualpanel}[2]{\placeholderpanel{manual/#1}{#2}}
\newcommand{\autopanel}[2]{\placeholderpanel{#1}{#2}}

\newcommand{\manualgraphic}[2]{%
    \IfFileExists{\manualevidencepath{#1}}{%
        \includegraphics[width=0.47\textwidth]{\manualevidencepath{#1}}%
    }{%
        \manualpanel{#1}{#2}%
    }%
}

\newcommand{\autographic}[2]{%
    \IfFileExists{\autoevidencepath{#1}}{%
        \includegraphics[width=0.47\textwidth]{\autoevidencepath{#1}}%
    }{%
        \autopanel{#1}{#2}%
    }%
}

\newcommand{\placeholderfigure}[2][0.88\textwidth]{%
\begin{figure}[H]
    \centering
    \fbox{\begin{minipage}[c][0.22\textheight][c]{#1}
        \centering \textbf{截图占位}\\[0.5em]
        \small 建议文件名：\path{#2}
    \end{minipage}}
\end{figure}
}

\newcommand{\evidencefigure}[3][0.95\textwidth]{%
\begin{figure}[H]
    \centering
    \subfigure[本人终端截图]{\manualgraphic{#2}{本人终端截图待补，可保留真实终端背景与上下文}}
    \hfill
    \subfigure[自动化证据截图]{\autographic{#2}{日志固化/自动化截图（可选）}}
    \caption{#3}
\end{figure}
}

\newenvironment{reportbox}[1]{%
    \par\smallskip\noindent\hrule\smallskip\noindent\textbf{#1}\par\smallskip
}{%
    \smallskip\noindent\hrule\par\smallskip
}

\newcommand{\codepath}[1]{\path{#1}}
\newcommand{\term}[1]{\textbf{#1}}
\newcommand{\cmdline}[1]{{\ttfamily\small\detokenize{#1}}}

\newenvironment{opsummary}{%
\begin{reportbox}{操作摘要（终端 / 位置 / 命令 / 等待点）}
\begingroup\small\sloppy
\begin{itemize}
    \setlength{\itemsep}{0.25em}
    \setlength{\topsep}{0.2em}
}{%
\end{itemize}
\endgroup
\end{reportbox}
}

\begin{document}

\cover

\begin{abstract}
本报告记录“实验二：ARP 缓存中毒攻击”的实验目标、环境搭建、实施步骤、关键现象与深入分析。实验重点包括：使用 ARP request / ARP reply / gratuitous ARP 三种方式投毒缓存；在双向投毒条件下，通过开关 IP forwarding 观察链路恢复与中断；并在 Telnet 与 Netcat 场景中实现中间人（MITM）流量篡改。

\textbf{关键词}: ARP, ARP poisoning, MITM, Scapy, IP forwarding
\end{abstract}

\tableofcontents
\newpage

\section{实验目的}
\begin{enumerate}
    \item 理解 ARP 的工作机制与“缓存更新”的关键触发条件；
    \item 通过构造 ARP request / reply / gratuitous ARP 实现 ARP 缓存中毒，并比较不同方式与不同缓存状态下的效果；
    \item 在双向投毒条件下验证 IP forwarding 对通信可达性的决定性作用；
    \item 在 MITM 场景中对 Telnet/Netcat 的 TCP payload 进行等长篡改，理解“长度不变”对 TCP 序列号一致性的意义；
    \item 从防御视角总结可检测点与缓解手段，为后续 TCP/DNS 等实验建立“前提与可见性”意识。
\end{enumerate}

\section{实验原理与背景}

\subsection{ARP 协议与缓存更新行为}
ARP（Address Resolution Protocol）用于在同一二层广播域内完成 \texttt{IP} $\rightarrow$ \texttt{MAC} 的地址解析。由于 ARP 协议本身缺少认证机制，主机通常会在收到 ARP request/reply 时\textbf{无条件或弱条件地更新缓存}，这使得攻击者能够通过伪造报文将“对方 IP”映射到“攻击者 MAC”。

\subsection{ARP 缓存中毒与 MITM}
一旦 A 的缓存中把 B 的 IP 映射到 M 的 MAC，则 A 发往 B 的以太网帧将首先到达 M。若 B 同样被投毒将 A 映射到 M，则双向流量都会经过 M，从而形成典型 MITM。此时：
\begin{itemize}
    \item 若 M 未开启 IP forwarding，则流量会在 M 处被内核丢弃，A/B 通信失败；
    \item 若 M 开启 IP forwarding，内核可以像“路由器”一样转发报文，A/B 通信恢复；
    \item 若 M 关闭 IP forwarding，但用户态脚本对报文进行“捕获+重放”，则可在转发过程中篡改 payload，实现可控 MITM。
\end{itemize}

\subsection{为什么篡改必须等长？}
TCP 的序列号/确认号严格依赖字节流偏移。若 MITM 篡改 payload 并改变长度，则从该包开始，后续所有方向上的 seq/ack 都会错位，除非攻击者对每个后续包持续修正序列号，否则连接会立即异常。因此本实验中对 Telnet（单字符）与 Netcat（整行）都采用“\textbf{等长替换}”策略，避免触发序列号漂移。

\section{实验环境与约束}

\subsection{容器拓扑}
实验使用 Docker Compose 构造三个容器（A/B/M）处于同一网段 \texttt{10.9.0.0/24}。其中 M 容器开启 \texttt{privileged} 并挂载共享目录 \codepath{lab2/Labsetup/volumes/} 以便在宿主机编辑 Scapy 脚本并在容器内直接运行。

\begin{table}[H]
\centering
\begin{tabularx}{\textwidth}{>{\raggedright\arraybackslash}p{0.20\textwidth} >{\raggedright\arraybackslash}p{0.22\textwidth} X}
\toprule
角色 & 容器名 / IP & 说明 \\
\midrule
Host A & \texttt{A-10.9.0.5} / \texttt{10.9.0.5} & 受害者 / 客户端（arp/ping/telnet/nc）\\
Host B & \texttt{B-10.9.0.6} / \texttt{10.9.0.6} & 受害者 / 服务器（telnet 服务、nc -l）\\
Host M & \texttt{M-10.9.0.105} / \texttt{10.9.0.105} & 攻击者（投毒 + MITM + sysctl）\\
\bottomrule
\end{tabularx}
\caption{实验容器与角色}
\end{table}

\subsection{终端与执行位置约定}
为保证步骤可复现并便于截图取景，本实验按终端编号组织操作：
\begin{table}[H]
\centering
\begin{tabularx}{\textwidth}{>{\raggedright\arraybackslash}p{0.10\textwidth} >{\raggedright\arraybackslash}p{0.30\textwidth} X}
\toprule
终端 & 执行位置 & 用途 \\
\midrule
\term{T0} & 宿主机 / WSL & \texttt{docker-compose}、抓包（可选）、编译报告 \\
\term{T1} & M（root） & 运行投毒脚本 \texttt{/volumes/arp/*} \\
\term{T1b} & M（root） & 并行运行 MITM 脚本（Telnet/Netcat） \\
\term{T2} & A & 查看 ARP 缓存、触发流量（ping/telnet/nc） \\
\term{T3} & B & 服务器端（nc -l、观测） \\
\bottomrule
\end{tabularx}
\caption{终端与执行位置摘要}
\end{table}

\subsection{证据与补图约定（简述）}
手工截图统一保存至 \codepath{lab2/evidence/manual/}，自动化证据（如日志渲染截图）保存至 \codepath{lab2/evidence/}。本文不在 PDF 内展开截图清单细则；以 \codepath{lab2/runbook_taskset1.md} 作为实际执行时的对照清单。

\section{实验结果总览}
\begin{table}[H]
\centering
\begin{tabularx}{\textwidth}{>{\raggedright\arraybackslash}p{0.13\textwidth} >{\raggedright\arraybackslash}p{0.20\textwidth} >{\raggedright\arraybackslash}p{0.14\textwidth} X}
\toprule
任务 & 目标 & 结果 & 关键观察 \\
\midrule
1.A & ARP request 投毒 & 成功 & A 会用 ARP request 的 sender 信息更新缓存（B $\rightarrow$ M）\\
1.B & ARP reply 投毒（两情景） & 部分成功 & 缓存已存在可覆盖；缓存缺失时不创建新条目（本环境观测）\\
1.C & Gratuitous ARP（两情景） & 部分成功 & 缓存已存在时可覆盖；缓存缺失时不创建新条目\\
2 & Telnet MITM（IP 转发开关） & 成功 & \texttt{ip\_forward=0} 时 ping 100\% 丢包；开启后恢复且 TTL $\downarrow 1$（需关 redirect）\\
3 & Netcat MITM（等长替换） & 成功 & 仅替换 \texttt{zhou} $\rightarrow$ \texttt{AAAA} 保持 TCP seq/ack 一致，连接稳定\\
\bottomrule
\end{tabularx}
\caption{实验二任务完成情况总览}
\end{table}

\section{任务 1：ARP 缓存中毒攻击}

\subsection{任务 1.A：使用 ARP request 投毒}
\subsubsection{实施步骤}
\begin{opsummary}
    \item \term{T2（A）}：先产生正常 ARP 条目并查看缓存\\
    \cmdline{ping -c 1 10.9.0.6}\\
    \cmdline{arp -n}
    \item \term{T1（M）}：发送伪造 ARP request\\
    \cmdline{cd /volumes/arp}\\
    \cmdline{python3 arp_poison.py --method request --victim-ip 10.9.0.5 --spoof-ip 10.9.0.6 --count 1 --show}
    \item \term{T2（A）}：再次查看缓存，观察 B 的 IP 是否指向 M 的 MAC\\
    \cmdline{arp -n}
\end{opsummary}

\subsubsection{结果与分析}
\evidencefigure{t2-1a-arp-request.png}{任务 1.A：使用 ARP request 投毒 A 的 ARP 缓存}

从实验现象可以看到：投毒前，A 的 ARP 缓存中 \texttt{10.9.0.6} 对应的是 B 的真实 MAC；投毒后，\texttt{10.9.0.6} 的 \texttt{HWaddress} 被替换为 M 的 MAC（与 \texttt{10.9.0.105} 的条目一致）。这说明在 Linux 的 ARP 实现中，\textbf{收到 ARP request 也会用 sender 信息（\texttt{psrc/hwsrc}）更新邻居缓存}。因此，即使我们发送的是“who-has”类型的 request，只要把 sender IP 伪装成 B，把 sender MAC 写成 M，就能诱导 A 把“B $\rightarrow$ M”的映射写入缓存。

这一点非常关键：ARP 的信任边界并不是“只信 reply”，而是“在很多实现中 request/reply 都会触发缓存学习”。从攻击角度看，request 方式往往比 reply 更“通用”，因为它不依赖受害者是否存在 pending ARP 请求。

\subsection{任务 1.B：使用 ARP reply 投毒（两种缓存情景）}

\subsubsection{情景 1：B 已在 A 的缓存中}
\begin{opsummary}
    \item \term{T2（A）}：确保缓存存在 B 条目\\
    \cmdline{ping -c 1 10.9.0.6}\\
    \cmdline{arp -n}
    \item \term{T1（M）}：发送伪造 ARP reply\\
    \cmdline{python3 arp_poison.py --method reply --victim-ip 10.9.0.5 --spoof-ip 10.9.0.6 --count 1 --show}
    \item \term{T2（A）}：查看缓存是否被替换\\
    \cmdline{arp -n}
\end{opsummary}
\evidencefigure{t2-1b-arp-reply-s1-cache-exists.png}{任务 1.B 情景 1：ARP reply 投毒（缓存已存在）}

实验结果显示：在“缓存已存在”的情况下，A 会接受该 ARP reply 并覆盖已有条目。\texttt{10.9.0.6} 从 B 的真实 MAC 更新为 M 的 MAC，说明该环境下已有邻居项可被伪造 reply 直接改写。

\subsubsection{情景 2：B 不在 A 的缓存中}
\begin{opsummary}
    \item \term{T2（A）}：删除条目并确认缓存中无 B\\
    \cmdline{arp -d 10.9.0.6}\\
    \cmdline{arp -n}
    \item \term{T1（M）}：发送伪造 ARP reply\\
    \cmdline{python3 arp_poison.py --method reply --victim-ip 10.9.0.5 --spoof-ip 10.9.0.6 --count 1 --show}
    \item \term{T2（A）}：观察系统是否创建/更新条目\\
    \cmdline{arp -n}
\end{opsummary}
\evidencefigure{t2-1b-arp-reply-s2-cache-miss.png}{任务 1.B 情景 2：ARP reply 投毒（缓存不存在）}

在“缓存不存在”的情况下，A 没有因为该 ARP reply 创建新条目，投毒失败。这说明本环境对 unsolicited ARP reply 的“凭空建项”仍较保守。结合两种情景可得：\textbf{ARP reply 在本环境能覆盖已有项，但不能在 cache miss 时直接新建项}，因此其可利用性强依赖受害者当前邻居状态。

这一现象的意义在于：ARP 投毒并不是“永远有效”的万能招数，它受到操作系统邻居子系统策略的直接影响；同一份攻击代码在不同系统/不同参数下可能得到完全不同的结果。报告中把这种差异写清楚，比“硬凑成功”更能体现对实验环境的理解。

\subsection{任务 1.C：使用 Gratuitous ARP 投毒（两种缓存情景）}

\subsubsection{情景 1：缓存已存在}
\begin{opsummary}
    \item \term{T2（A）}：确保缓存存在 B 条目\\
    \cmdline{ping -c 1 10.9.0.6}\\
    \cmdline{arp -n}
    \item \term{T1（M）}：发送 gratuitous ARP\\
    \cmdline{python3 arp_poison.py --method gratuitous --victim-ip 10.9.0.5 --spoof-ip 10.9.0.6 --count 1 --show}
    \item \term{T2（A）}：查看缓存是否变化\\
    \cmdline{arp -n}
\end{opsummary}
\evidencefigure{t2-1c-gratuitous-s1-cache-exists.png}{任务 1.C 情景 1：Gratuitous ARP 投毒（缓存已存在）}

在情景 1 中，gratuitous ARP 成功覆盖了 A 的已有条目：\texttt{10.9.0.6} 被更新为 M 的 MAC。可将其理解为“对已有映射的公告更新”：即使 \texttt{op} 仍是 request（who-has），但由于 \texttt{psrc=pdst} 并广播发送，接收方会把它当作一种地址变更通告来处理。

\subsubsection{情景 2：缓存不存在}
\begin{opsummary}
    \item \term{T2（A）}：删除条目并确认\\
    \cmdline{arp -d 10.9.0.6}\\
    \cmdline{arp -n}
    \item \term{T1（M）}：发送 gratuitous ARP\\
    \cmdline{python3 arp_poison.py --method gratuitous --victim-ip 10.9.0.5 --spoof-ip 10.9.0.6 --count 1 --show}
    \item \term{T2（A）}：观察是否创建/更新条目\\
    \cmdline{arp -n}
\end{opsummary}
\evidencefigure{t2-1c-gratuitous-s2-cache-miss.png}{任务 1.C 情景 2：Gratuitous ARP 投毒（缓存不存在）}

在情景 2 中，当 A 的缓存里根本没有 \texttt{10.9.0.6} 条目时，gratuitous ARP 并未创建新条目，投毒失败。也就是说，在本环境里 gratuitous ARP 更像是“覆盖已有项”的更新机制，而不是“无条件新增项”的学习机制。

因此，从攻击链角度看：若希望利用 gratuitous ARP 稳定投毒，通常需要先触发受害者与目标主机有过一次正常通信（例如先 ping 一次），让缓存里出现该条目，再用 gratuitous ARP 覆盖；或者直接采用任务 1.A 的 ARP request 投毒（本环境中可在无缓存时直接创建条目），在工程上更稳健。

\section{任务 2：使用 ARP 缓存中毒在 Telnet 实施 MITM}

\subsection{双向投毒与 IP 转发开关实验}
\begin{opsummary}
    \item \term{T1（M）}：持续双向投毒（保持运行）\\
    \cmdline{cd /volumes/arp}\\
    \cmdline{python3 arp_mitm_poison.py --method request --interval 1 --restore}
    \item \term{T1（M）}：关闭 IP forwarding 并测试\\
    \cmdline{sysctl net.ipv4.ip_forward=0}
    \item \term{T2（A）}：ping B\\
    \cmdline{ping -c 2 -W 1 10.9.0.6}
    \item \term{T1（M）}：开启 IP forwarding 并重复测试\\
    \cmdline{sysctl net.ipv4.ip_forward=1}
\end{opsummary}
\evidencefigure{t2-2-ping-ipfwd0.png}{任务 2：双向投毒下关闭 IP forwarding 的 ping 现象}
\evidencefigure{t2-2-ping-ipfwd1.png}{任务 2：双向投毒下开启 IP forwarding 的 ping 现象}

从结果可见，当双向投毒生效后，A 与 B 的以太网帧目的 MAC 都会指向 M。此时：
\begin{itemize}
    \item \textbf{\texttt{ip\_forward=0}}：M 仍会接收到目的 IP 为对方主机的 IP 包，但内核不进行转发，因此 ping 直接变为 100\% 丢包。该现象说明 ARP 投毒不仅可用于窃听，也可作为“链路层拒绝服务（DoS）”手段。
    \item \textbf{\texttt{ip\_forward=1}}：M 的内核开始转发数据包，A 与 B 通信恢复。由于 IP 包经过了一次转发，TTL 会在 M 处减 1，因此我们在 ping 输出中观察到 TTL 由 64 变为 63，这是“流量确实绕行 M”的直接证据。
\end{itemize}

需要额外记录的一点是：Linux 在“同网段转发”时可能发送 ICMP redirect，提示主机“其实可以直连对方”。为了避免干扰实验现象并降低 MITM 暴露度，本实验额外设置 \texttt{send\_redirects=0} 后再进行开启转发的对照测试。

\subsection{实施 Telnet MITM：把输入替换为 Z}
\begin{opsummary}
    \item \term{T1（M）}：保持投毒运行；先开启转发用于建连\\
    \cmdline{sysctl net.ipv4.ip_forward=1}
    \item \term{T2（A）}：建立 Telnet 并登录\\
    \cmdline{telnet 10.9.0.6}\\
    \cmdline{user: seed / pass: dees}
    \item \term{T1（M）}：连接建立后关闭转发\\
    \cmdline{sysctl net.ipv4.ip_forward=0}
    \item \term{T1b（M）}：启动用户态 MITM 程序（捕获并转发 TCP，同时篡改 A->B payload）\\
    \cmdline{python3 mitm_telnet.py --show}
    \item \term{等待/停止}：在 Telnet 内随便输入字符观察是否变为 Z；完成截图后 Ctrl+C 退出
\end{opsummary}
\evidencefigure{t2-2-telnet-mitm.png}{任务 2：Telnet MITM（客户端输入被替换为 Z）}

该步骤的关键在于“\textbf{关内核转发、开用户态转发}”。一旦 Telnet 连接建立，我们关闭 \texttt{ip\_forward}，使得内核不再替我们转发 A/B 的 TCP 包；然后由 \texttt{mitm\_telnet.py} 在用户态捕获并转发报文，并在 A$\rightarrow$B 方向将可打印字符逐字节替换为 \texttt{Z}（长度不变）。因此，A 在终端中输入 \texttt{abc}，最终显示为 \texttt{ZZZ}，且服务器端实际接收到的也是 \texttt{ZZZ}，从而达到“输入被篡改但连接仍稳定”的效果。

\section{任务 3：使用 ARP 缓存中毒在 Netcat 实施 MITM}

\begin{opsummary}
    \item \term{T1（M）}：保持双向投毒运行；先开启转发用于建连\\
    \cmdline{sysctl net.ipv4.ip_forward=1}
    \item \term{T3（B）}：启动 netcat 服务端\\
    \cmdline{nc -lp 9090}
    \item \term{T2（A）}：连接到 B\\
    \cmdline{nc 10.9.0.6 9090}
    \item \term{T1（M）}：连接建立后关闭转发\\
    \cmdline{sysctl net.ipv4.ip_forward=0}
    \item \term{T1b（M）}：启动 MITM 程序（仅等长替换 \texttt{zhou} $\rightarrow$ \texttt{AAAA}）\\
    \cmdline{python3 mitm_netcat.py --surname zhou --show}
    \item \term{等待/停止}：在 A 中输入包含 \texttt{zhou} 的一行文本并观察 B 输出；完成截图后 Ctrl+C 退出
\end{opsummary}
\evidencefigure{t2-3-netcat-mitm.png}{任务 3：Netcat MITM（姓氏拼音被替换为等长 A 串）}

与 Telnet 类似，Netcat 场景本质仍是“捕获$\rightarrow$修改$\rightarrow$重放”。差异在于：Netcat 往往按“行”发送数据，因此本任务用正则在 payload 中定位 \texttt{zhou} 并替换为 \texttt{AAAA}（严格等长）。只要长度一致，TCP 序列号就无需额外修正，连接可以持续工作，这也解释了手册强调“长度必须与你姓拼音长度相同”的原因。

\section{实验分析与讨论（要点提纲）}
\begin{itemize}
    \item \textbf{三种投毒方式对比（结合本环境）：} request 方式最稳定，既能覆盖已有条目，也能在无条目时诱导创建；reply 方式表现为“可覆盖已有项、但 cache miss 不建项”；gratuitous 介于两者之间，可覆盖已有项，但对“无缓存新增”不生效。
    \item \textbf{IP forwarding 的作用：} 双向投毒后，M 接收到的 IP 包目的地址并非自身，只有在 \texttt{ip\_forward=1} 时内核才会继续转发；否则通信立刻中断，这也是 MITM 能“卡断链路”的关键。
    \item \textbf{ICMP redirect 的工程细节：} 当 M 在同网段转发时，默认可能发送 redirect，提示主机走“更优路径”。关闭 \texttt{send\_redirects} 有助于保持实验现象稳定，也更符合真实 MITM 的隐蔽需求。
    \item \textbf{用户态转发与等长篡改：} Telnet/Netcat 的 MITM 本质是“捕获 $\rightarrow$ 修改 $\rightarrow$ 重新注入”。只要 payload 长度不变，TCP 序列号无需修正，连接就能稳定维持；若长度变化，必须实现持续的 seq/ack 重写。
    \item \textbf{可检测点与防御：} ARP 表异常抖动、同一 IP 对应 MAC 频繁变化、网关/主机启用动态 ARP 检测（DAI）、静态 ARP 绑定、交换机侧 DHCP snooping + DAI 等，均可缓解或检测此类攻击。
\end{itemize}

\section{补充实验设计与扩展验证（可选）}
\begin{enumerate}
    \item \textbf{投毒持久性实验：} 在不持续投毒的情况下测量 ARP 缓存条目回收/刷新时间；比较不同系统/配置差异。
    \item \textbf{非等长篡改对照：} 故意让 payload 变长/变短，观察 TCP 连接如何快速失效，并分析 seq/ack 错位导致的现象。
    \item \textbf{检测对照：} 在宿主机桥接接口抓包，统计 ARP reply 的频率与来源，尝试用简单规则识别“异常 ARP 更新”。
\end{enumerate}

\section{思考题与深入回答}

\subsection{为什么 ARP 天生脆弱？}
ARP 的脆弱性来自其设计目标：它要在二层广播域里快速完成地址解析，因此协议本身既没有认证，也没有加密。攻击者只要处于同一广播域，就能构造任意 \texttt{psrc/hwsrc} 组合的 ARP 报文，并试图触发受害者缓存更新。对于接收方而言，它无法通过 ARP 报文本身验证“某个 IP 是否真的属于某个 MAC”，因为这一绑定关系并不存在可被查询的权威来源；交换机只负责转发帧，不会为 ARP 绑定背书。因此，\textbf{ARP 报文中的 sender 字段几乎都可以被伪造}，而接收方最多只能用“是否曾请求过/是否已有条目/是否允许覆盖”等策略做弱防御。

\subsection{为什么同样是 ARP reply，不同缓存状态下可能出现不同结果？}
是否接受 ARP reply 更新缓存，取决于操作系统的邻居子系统策略以及相关 sysctl 配置。常见差异包括：
\begin{itemize}
    \item \textbf{是否接受 unsolicited reply：} 若 \texttt{arp\_accept=0}，系统可能直接忽略没有触发请求的 ARP reply；
    \item \textbf{是否允许覆盖已存在条目：} 有的实现会用新映射覆盖旧映射，有的实现会保持原映射不变或进入“可疑/延迟更新”流程；
    \item \textbf{条目状态机差异：} Linux 邻居表有 \texttt{REACHABLE/STALE/INCOMPLETE} 等状态，不同状态下对更新的处理不同。
\end{itemize}
因此，“是否成功”并不是 ARP reply 格式对不对，而是受害系统愿不愿意学习/覆盖。在本实验环境中，我们观察到：缓存已存在时可被 ARP reply 覆盖；缓存缺失时不会凭空建项，体现了邻居状态机与策略控制共同决定投毒效果。

\subsection{为什么实验中强调“等长替换”？如果必须变长，MITM 需要额外维护哪些状态？}
TCP 序列号表示字节流偏移。若把 payload \textbf{变长/变短}，从该包开始两端看到的“已发送字节数”就不同步，后续 ACK 会持续错位，连接通常立即进入重传或直接断开。

若必须变长，MITM 至少需要：
\begin{enumerate}
    \item 维护每个方向的“累计偏移差”（delta），并对后续所有包持续修正 seq/ack；
    \item 处理乱序、重传与窗口更新等情况，保证修正逻辑对重复包一致；
    \item 在应用层边界（如一行/一条消息）处设计更复杂的同步策略，避免拆包/粘包导致替换跨包。
\end{enumerate}
这也是为什么本实验选择等长替换：把难点集中在“能否稳定捕获并重放正确的 TCP 包”，而不引入完整 TCP 代理所需的大量状态维护。

\section{实验总结}
本实验通过“投毒 $\rightarrow$ 链路重定向 $\rightarrow$ 转发开关 $\rightarrow$ 用户态篡改”的完整链路，验证了 ARP 无认证导致的根本性风险，并将 Lab 1 的“嗅探与伪造”能力提升为可控的 MITM 攻击能力。更重要的是，本实验让我们明确：很多上层攻击（TCP/DNS 等）并不是孤立发生，而是依赖于链路层/路由层前提先被攻击者控制或劫持。

\appendix

\section{附录：关键脚本路径}
\begin{lstlisting}[language=bash, caption={实验二关键脚本（共享目录）}]
# 攻击者容器 M 内
cd /volumes/arp
ls -la

# Task 1: 单次投毒
python3 arp_poison.py --method request --victim-ip 10.9.0.5 --spoof-ip 10.9.0.6

# MITM: 双向投毒
python3 arp_mitm_poison.py --method request --interval 1 --restore

# Task 2/3: 用户态 MITM
python3 mitm_telnet.py --show
python3 mitm_netcat.py --surname zhou --show
\end{lstlisting}

\section{附录：手工复现清单}
面向人的逐步命令与截图清单见 \codepath{lab2/runbook_taskset1.md}。

\end{document}
